% Generated draft: requires mathematical and source review.
\documentclass[11pt]{article}
\usepackage[margin=25mm]{geometry}
\usepackage{fontspec}
\setmainfont{Latin Modern Roman}
\usepackage{amsmath,amssymb,mathtools,unicode-math}
\setmathfont{Latin Modern Math}
\usepackage{enumitem,needspace,xurl,hyperref}
\setlength{\parindent}{0pt}
\setlength{\parskip}{5pt}
\setlength{\emergencystretch}{3em}
\newcommand{\N}{\mathbb{N}}
\newcommand{\Z}{\mathbb{Z}}
\newcommand{\Q}{\mathbb{Q}}
\newcommand{\R}{\mathbb{R}}
\newcommand{\C}{\mathbb{C}}
\newcommand{\F}{\mathbb{F}}
\newcommand{\A}{\mathbb{A}}
\newcommand{\PP}{\mathbb{P}}
\newcommand{\E}{\mathbb{E}}
\newcommand{\eps}{\varepsilon}
\newcommand{\dd}{\,\mathrm{d}}
\begin{document}
\section*{UnsolvedMath ID: 700815}
\textbf{Source problem code:} catalogue entry
\subsection{Definitions and mathematical statement}
For a group \(H\), let \(\operatorname{Aut}(H)\) denote the group of all bijective group homomorphisms \(H\to H\), with composition as its operation. Given a group \(G\), define groups recursively by
\[
A_1(G)=G,\qquad A_{i+1}(G)=\operatorname{Aut}\bigl(A_i(G)\bigr)\quad(i\in\Z_{\geq 1}).
\]
The problem asks whether there exists a finite group \(G\) for which the set of isomorphism classes
\[
\bigl\{[A_i(G)]_{\cong}:i\in\Z_{\geq 1}\bigr\}
\]
is infinite. Equivalently, does there exist a finite group \(G\) and a strictly increasing sequence \(n_1<n_2<\cdots\) of positive integers such that
\[
A_{n_r}(G)\not\cong A_{n_s}(G)\qquad\text{for all distinct }r,s\in\Z_{\geq 1}?
\]
Here \([H]_{\cong}\) denotes the isomorphism class of the abstract group \(H\); the recursion specifies groups and does not require connecting embeddings.
\subsection{Short English statement}
Can a finite group yield infinitely many distinct group isomorphism classes by repeatedly taking its automorphism group? The supplied record calls this a well-known problem, attributes its proposal to M.~Short in 1990, and labels it open; that status is supplied source information, not an independently verified assessment.
\subsection{Sources}
\begin{itemize}
\item Local UnsolvedMath registry, supplied complete source record: permanent ID \(700815\), exact source problem code \texttt{catalogue entry}, title \emph{Kourovka Notebook Problem 11.123}. The record supplies the status \texttt{open}.
\item \emph{The Kourovka Notebook (Unsolved Problems in Group Theory)}, 11th issue (1990), Problem 11.123; proposed by M.~Short and described in the supplied statement as a ``Well-known problem.''
\item The supplied background also identifies the 21st issue edition (2026), \texttt{arXiv:1401.0300v48}. The collection description names E.~I.~Khukhro and V.~D.~Mazurov as editors of the Issues 1--21 collection (1965--2026). These bibliographic details are supplied attribution and have not been independently checked.
\end{itemize}
\end{document}
